% ia-utilisation.tex — Une page : comment l'IA a été utilisée pour développer le
% site de simulations (outil, cadre, cycle de travail, vérifications, rôle de
% l'enseignant, limites). Sert aussi de déclaration d'utilisation de l'IA.
% Compilation : pdflatex ia-utilisation.tex (deux fois), ou
%   powershell -ExecutionPolicy Bypass -File memoire\ia-utilisation\compiler.ps1

\documentclass[10pt, a4paper]{article}

\usepackage[T1]{fontenc}
\usepackage[utf8]{inputenc}
\usepackage{lmodern}
\usepackage[french]{babel}
\usepackage{microtype}
\usepackage[a4paper, margin=1.8cm]{geometry}
\usepackage{enumitem}
\usepackage{xcolor}
\usepackage{csquotes}
\usepackage[hidelinks]{hyperref}

\setlist{itemsep=0.15em, topsep=0.25em}
\setlength{\parskip}{0.3em}
\setlength{\parindent}{0pt}
\pagestyle{empty}

\begin{document}

{\Large\bfseries Utilisation de l'intelligence artificielle dans le projet\par}
\vspace{0.4em}

\textbf{Outil et principe.} Le code du site a été écrit avec Claude Code
(Anthropic), un assistant de programmation fondé sur un modèle de langage,
utilisé dans l'éditeur VS Code. L'enseignant décrit en français ce qu'il veut
(un exercice, une correction, une présentation), l'assistant lit le dépôt,
propose ou écrit le code, lance les tests, puis rend compte de ce qu'il a fait.
L'enjeu n'est pas de \enquote{faire écrire du code} à l'IA, mais d'encadrer son
travail pour que le résultat soit juste en physique, vérifiable et adapté à des
étudiants de BAC1.

\textbf{Un cadre écrit avant le code.} Les règles du projet sont fixées dans
des fichiers lus à chaque session (\texttt{CLAUDE.md}) :
\begin{itemize}
  \item \emph{pédagogie} : public BAC1, notations du formulaire du cours
        EPL1201, formules affichées sous trois formes (symbole, valeurs
        substituées, résultat), hypothèses écrites sur la page ;
  \item \emph{technique} : HTML, CSS et JavaScript sans dépendance, pages
        utilisables en double-cliquant (\texttt{file://}), calculs physiques
        séparés de l'affichage et jamais partagés entre exercices ;
  \item \emph{méthode} : tests et vérifications confiés à des agents séparés,
        aucun commit sans accord, et toute correction de plus d'une vingtaine
        de lignes expliquée avant d'être faite.
\end{itemize}
Des procédures réutilisables (\enquote{skills}) décrivent les tâches
récurrentes : créer un exercice, vérifier la physique, écrire les tests
d'intégration, ajouter un mode jeu, contrôler les critères de fin
(\emph{definition of done}).

\textbf{Le cycle de travail.} Chaque demande suit le même cycle :
\begin{enumerate}
  \item l'enseignant formule la demande et tranche les choix pédagogiques (par
        exemple : un seau d'eau plutôt qu'une bille ; une boucle en clothoïde
        ou en cercle) ;
  \item l'assistant écrit le code, en mesurant ses propositions par des
        simulations quand le comportement physique est en jeu ;
  \item un agent distinct (le même modèle, lancé dans un contexte séparé, sans
        l'historique de la session) vérifie le code :
        tests unitaires (183 tests, \texttt{node -{}-test}), tests d'intégration
        dans un navigateur sur 9 profils d'appareils (Playwright), et un agent
        de vérification de la physique qui redérive les formules et les
        recoupe numériquement ;
  \item l'enseignant teste la page, relit, demande des corrections, et décide
        seul du commit.
\end{enumerate}

\textbf{Ce que les vérifications ont trouvé.} La séparation entre l'agent qui
code et celui qui vérifie a fait apparaître de vraies erreurs. Exemples sur le
looping avec wagon : une carte de formule de la clothoïde affichait une
réaction au sommet jusqu'à 22 fois trop grande ; dans le jeu du seau, le
frottement d'une articulation exerçait sur le seau une force qui ne passait
pas par la corde, contraire à la règle physique fixée par l'enseignant ; une
idée de départ \enquote{brutal}, testée à la demande de l'enseignant, faisait
tomber l'eau quel que soit le réglage et a été abandonnée sur la base des
mesures.

\textbf{Ce qui reste humain.} Les objectifs pédagogiques, le choix des
exercices, des notations et des hypothèses, l'acceptation de chaque résultat
et la responsabilité de son exactitude restent ceux de l'enseignant. L'IA
accélère l'écriture et la vérification ; elle ne remplace ni la compréhension
de la physique, ni la relecture.

\textbf{Limites.} L'assistant se trompe : certaines propositions ont échoué
(une commande de moteur à vitesse constante rendait le jeu instable) et
n'ont été écartées que parce qu'elles étaient mesurées. Les vérifications
réduisent le risque sans l'annuler, et un code qu'aucun humain ne relit reste
un code non vérifié.

\textbf{Données.} Seul le code du site et ses documents sont transmis au
fournisseur du modèle ; aucune donnée d'étudiant n'est utilisée.

\textbf{Déclaration.} Le code du site, ce texte et le rapport technique sur le
jeu du seau d'eau ont été rédigés avec Claude Code (Anthropic, modèles Claude,
de mai à novembre 2026). L'auteur a fixé le cadre, relu les formules, les tests
et le comportement des pages, et en assume l'exactitude. [À COMPLÉTER :
vérifier la forme de déclaration exigée par le promoteur ou la faculté.]

\end{document}
